\documentclass[11pt]{article}
\usepackage[T1]{fontenc}
\usepackage[margin=1.8cm]{geometry}
\usepackage{amsmath,amssymb}
\setlength{\parindent}{0pt}
\setlength{\parskip}{8pt}
\begin{document}

\begin{center}
{\Large\bfseries GEOMETR\'IA 8 -- CLAVE DOCENTE}\\[3pt]
{\large \'Areas sombreadas, escala, \'area y per\'imetro integrado}
\end{center}

\section*{Semana 4 -- \'Areas sombreadas simples}
\textbf{1.} $A=10^2-\pi(3)^2=100-9\pi\approx71.73\,\mathrm{cm}^2$.

\textbf{2.} $A=15(8)-\frac12\pi(4)^2=120-8\pi\approx94.87\,\mathrm{cm}^2$.

\textbf{3.} $A=12^2-\frac14\pi(12)^2=144-36\pi\approx30.90\,\mathrm{cm}^2$.

\textbf{4.} $A=\pi(7)^2-8^2=49\pi-64\approx89.94\,\mathrm{cm}^2$.

\section*{Semana 5 -- \'Areas sombreadas complejas}
\textbf{5.} $A=\pi(10^2-6^2)=64\pi\approx201.06\,\mathrm{cm}^2$.

\textbf{6.} Cuatro cuadrantes de radio $8$ forman un c\'irculo: $A=16^2-\pi(8)^2=256-64\pi\approx54.94\,\mathrm{cm}^2$.

\textbf{7.} $A=18(10)-2\pi(2)^2=180-8\pi\approx154.87\,\mathrm{cm}^2$.

\textbf{8.} $A=\pi(8)^2-2\pi(2)^2=56\pi\approx175.93\,\mathrm{cm}^2$.

\textbf{9.} Dos semic\'irculos de radio $6$ equivalen a un c\'irculo: $A=24(12)-\pi(6)^2=288-36\pi\approx174.90\,\mathrm{cm}^2$.

\section*{Semana 6 -- Per\'imetro y \'area bajo escala}
\textbf{10.} $k=15/5=3$. El per\'imetro cambia por $\times3$ y el \'area por $\times3^2=\times9$.

\textbf{11.} $k=2/4=0.5$. La circunferencia cambia por $\times0.5$ y el \'area por $\times(0.5)^2=\times0.25$.

\textbf{12.} $A'=k^2A=(1.5)^2(40)=90\,\mathrm{cm}^2$.

\textbf{13.} $P'=2.5(36)=90\,\mathrm{cm}$; $A'=(2.5)^2(50)=312.5\,\mathrm{cm}^2$.

\section*{Semana 7 -- Reto integrado}
\textbf{14.} La pista es un rect\'angulo $24\times10$ m\'as un c\'irculo de radio $5$: $A=240+25\pi\approx318.54\,\mathrm{m}^2$. El per\'imetro es $P=2(24)+2\pi(5)=48+10\pi\approx79.42\,\mathrm{m}$.

\textbf{15.} Superficie libre: $A=18(14)-\pi(4)^2=252-16\pi\approx201.73\,\mathrm{m}^2$. Per\'imetro exterior: $P=2(18+14)=64\,\mathrm{m}$.

\textbf{16.} Con $k=2.5$, el per\'imetro cambia por $\times2.5$ y el \'area por $\times6.25$.

\textbf{17.} \'Area de la corona: $A=\pi(9^2-4^2)=65\pi\approx204.20\,\mathrm{m}^2$. Longitud total de ambos bordes: $L=2\pi(9)+2\pi(4)=26\pi\approx81.68\,\mathrm{m}$.

\section*{Exit ticket -- respuestas esperadas}
\textbf{a)} Borde, contorno, cerca, cinta alrededor, longitud del l\'imite.\\
\textbf{b)} Superficie, cubrir, pintar, piso, regi\'on, \'area sombreada.\\
\textbf{c)} Si $k=4$, el \'area cambia por $k^2=16$.

\end{document}
